\documentclass[12pt,a4paper]{exam}
\usepackage[utf8]{inputenc}
\usepackage[T1]{fontenc}
\usepackage[french]{babel}
\usepackage{amsmath,amssymb}
\usepackage{multicol}
\usepackage{geometry}
\usepackage{pgfplots}
\usepackage{cancel}
\usepackage{xcolor}

\pgfplotsset{compat=1.18}
\geometry{margin=2cm}

\header{}{Interrogation de mathématiques}{2nde8}
\footer{}{\thepage}{\it{G.M.}}

\begin{document}

\begin{center}
    \Large\textbf{Correction interrogation\\Équations et inéquations}
\end{center}

\vspace{0.3cm}

\noindent
\textbf{Durée :} 1h \hfill
\textbf{Calculatrice :} \underline{NON} autorisée

\vspace{0.5cm}

% ======================================================
\underline{Remarques générales :}
\vspace{5pt}
{\color{red}
\begin{itemize}
    \item Pensez à simplifier vos fractions \textbf{dès que c'est possible}. Par exemple dans 
    l'exercice 1 question b) j'ai eu beaucoup de $\dfrac{10}{6}$. Ce 
    n'est pas simplifié au maximum.
    \item Je ne mets pas les tableaux de signes dans la correction, à vous de les refaire pour retrouver 
    les bons résultats
\end{itemize}
}

\vspace{0.5cm}

% ======================================================
\noindent
\textbf{Exercice 1 : Équations (\ \ \ \  / 10 points )}
\vspace{5pt}

\begin{questions}

\question[1] Résoudre les équations suivantes dans $\mathbb{R}$ :
\begin{multicols}{2}
    \begin{parts}
        \part $x + 3 = 3x$
        \textcolor{red}{$S = \left\{\dfrac{3}{2}\right\}$}

        \part $\dfrac{6}{5}x - 2 = 0$
        \textcolor{red}{$S = \left\{\dfrac{5}{3}\right\}$}
    \end{parts}
\end{multicols}
\begin{multicols}{2}
    \begin{parts}
        \setcounter{partno}{2}
        \part $x^2 - 5 = 0$
        \textcolor{red}{$S = \{-\sqrt{5}; \sqrt{5}\}$}

        \part $(2x + 2)(3x - 4) = 0$
        \textcolor{red}{$S = \{-1; \dfrac{4}{3}\}$}
    \end{parts}
\end{multicols}
\begin{multicols}{2}
    \begin{parts}
        \setcounter{partno}{4}
        \part $\dfrac{2x - 1}{x+8} = 0$
        \textcolor{red}{$S = \left\{\dfrac{1}{2}\right\}$ avec $x \neq -8$}

        \part $\dfrac{x(3x+2)}{2x-3} = 0$
        \textcolor{red}{$S = \left\{0; -\dfrac{2}{3}\right\}$ avec $x \neq \dfrac{3}{2}$}
    \end{parts}
\end{multicols}

\end{questions}

\vspace{0.5cm}

% ======================================================
\noindent
\textbf{Exercice 2 : Inéquations (\ \ \ \  / 10 points )}
\vspace{5pt}

\begin{questions}

\question[4] Résoudre les inéquations suivantes dans $\mathbb{R}$ et indiquer les solutions sous forme d'intervalles : $S = ...$
\begin{multicols}{2}
    \begin{parts}
        \part $6x + 3 > 2x - 5$
        \textcolor{red}{$S = \left]-\dfrac{8}{4}; +\infty\right[$}

        \part $-3x \geq 3x - 1$
        \textcolor{red}{$S = \left]-\infty; \dfrac{1}{6}\right]$}
    \end{parts}
\end{multicols}
\begin{parts}
    \setcounter{partno}{2}
    \part $(1-2x)(4x+2) \leq 0$
    \textcolor{red}{$S = \left]-\infty; -\dfrac{1}{2}\right] \cup \left[\dfrac{1}{2}; +\infty\right[$}

    \part $\dfrac{2x}{-8x+16} < 0$
    \textcolor{red}{$S = \left]0; 2\right[$ avec $x \neq 2$}
\end{parts}
\begin{parts}
    \setcounter{partno}{4}
    \part $\dfrac{-x(4x+1)}{x^2+2} > 0$
    \textcolor{red}{$S = \left]-\dfrac{1}{4}; 0\right[$}

    \part $(x+2)(2x-3)(x-6) \geq 0$
    \textcolor{red}{$S = \left[-2; \dfrac{3}{2}\right] \cup \left[6; +\infty\right[$}
\end{parts}

\end{questions}

\vspace{10pt}

% ======================================================
\noindent
\textbf{Exercice 3 : BONUS (\ \ \ \  / 2 points )}
\vspace{5pt}

\begin{questions}

\question[1] On considère la fonction $g$ suivante :
\[
g(x) = \dfrac{1}{x} - \dfrac{2}{5-x} \quad \text{pour tout $x \neq 0$ et $x \neq 5$}
\]
Résoudre $g(x) \leq 0$.
\textcolor{red}{
\vspace{5pt}
\\
Il faut mettre sous le même dénominateur pour n'avoir qu'une seule fraction. On obtient :
\[
g(x) = \dfrac{(5-x) - 2x}{x(5-x)} = \dfrac{5-3x}{x(5-x)}
\]
On résout ensuite l'inéquation $\dfrac{5-3x}{x(5-x)} \leq 0$ en utilisant un tableau de signes.
\\$S = \left]-\infty; 0\right[ \cup \left[\dfrac{5}{3}; 5\right[$
}

\vspace{5pt}

\question[1] Résoudre les équations suivantes dans $\mathbb{R}$ :
\begin{multicols}{2}
    \begin{parts}
        \part $-3x^2+18x-27=0$
        \textcolor{red}{
        \\
        \vspace{5pt}On factorise par $-3$ : $-3(x^2 - 6x + 9) = 0$ puis on reconnaît une identité remarquable.
        \\
        $S = \{3\}$}

        \part $(11x-3)^2 = 11$
        \textcolor{red}{
        \\
        \vspace{5pt}On écrit $(11x-3)^2 - 11 = 0$ et on reconnaît une différence de carrés.
        \\
        $S = \left\{\dfrac{3-\sqrt{11}}{11}; \dfrac{3+\sqrt{11}}{11}\right\}$}
    \end{parts}
\end{multicols}

\end{questions}

\vspace{10pt}

\begin{center}
    \textit{Fin du devoir}
\end{center}

\end{document}
